%% Part VI -- Outlook (comic pages 44 to 49).
%% The section that stops describing and starts demanding: a device that
%% augments rather than entertains, women in the room, AI read as an economic
%% form rather than a technology, the ART framework, and the literacy without
%% which none of the rest can be argued about.
%%
%% ONE PANEL CAPTION, ONE SLIDE. The comic sets every topic page as six
%% panels in a 2x3 grid with one caption under each. Each caption takes a
%% slide of its own so the room reads one claim at a time; the six-dot row
%% under it marks which panel is on screen, and six consecutive slides
%% sharing a frame title are one comic page, read in the comic's own order.
%% Captions are verbatim, checked against a render of each page.
%% \setpagethesis opens each run with that page's one line, mine, not theirs.
%% Fragment only: no preamble, \input by ai_for_all.tex.

\sectionsubtitle{pages 44 to 49}
\section{Outlook}

% The budget follows the purchase decision and the purchase decision follows
% entertainment, so the device we got is the device we bought. Panel 4 is a
% product-management claim.
\setpagethesis{Your phone analyses you all day. What has it ever helped you think?}
\begin{frame}{Overall Human-Machine Intelligence}
  \panel{page 45, the smartphone as marketplace}
  \panelslide{45}{1}{We mainly interact with AI via our smartphone. A proprietary marketplace for countless uncurated applications (=platform economy).}
\end{frame}
\begin{frame}{Overall Human-Machine Intelligence}
  \panel{page 45, the smartphone as marketplace}
  \panelslide{45}{2}{On the smartphone, we use AI basically for entertainment (image recognition), communication (language and translation) and information (e.g. google).}
\end{frame}
\begin{frame}{Overall Human-Machine Intelligence}
  \panel{page 45, the smartphone as marketplace}
  \panelslide{45}{3}{Unbeknownst to us, AI improves performance and battery life, too. And: It analyzes us, its users.}
\end{frame}
\begin{frame}{Overall Human-Machine Intelligence}
  \panel{page 45, the smartphone as marketplace}
  \panelslide{45}{4}{Up to date, there is no smartphone that is especially good at helping us get better ideas or interpret information. There is no interface for that either.}
\end{frame}
\begin{frame}{Overall Human-Machine Intelligence}
  \panel{page 45, the smartphone as marketplace}
  \panelslide{45}{5}{The development budgets are driven by numbers of purchase. In turn, our purchase decision is based on entertainment and communication values, and the design of the device.}
\end{frame}
\begin{frame}{Overall Human-Machine Intelligence}
  \panel{page 45, the smartphone as marketplace}
  \panelslide{45}{6}{But couldn't we also develop a smartphone that makes us more creative, more intelligent? That understands like us? Really augment us?}
\end{frame}

% Opens on the human computers who ran the first calculations and closes on
% what a homogeneous development team reproduces. The parenthesis in panel 6 is
% the authors' own.
\setpagethesis{The first computers were women. The fear is not a takeover but a homogeneous room.}
\begin{frame}{Feminism}
  \panel{page 46, Ada Lovelace and Grace Hopper}
  \panelslide{46}{1}{The contribution of women to AI is usually underestimated, especially concerning the democratic applicability of computing.}
\end{frame}
\begin{frame}{Feminism}
  \panel{page 46, Ada Lovelace and Grace Hopper}
  \panelslide{46}{2}{They were the first computers, actually: running calculations in the first information network of the world.}
\end{frame}
\begin{frame}{Feminism}
  \panel{page 46, Ada Lovelace and Grace Hopper}
  \panelslide{46}{3}{What we fear is not an AI-takeover. We fear scientists' homogeneity leading them to (finally!) `create' intelligent beings -- in an egocentric way.}
\end{frame}
\begin{frame}{Feminism}
  \panel{page 46, Ada Lovelace and Grace Hopper}
  \panelslide{46}{4}{In an AI world, where `female' service robots are `hot' and submissive, how can children learn to treat women with respect, with dignity?}
\end{frame}
\begin{frame}{Feminism}
  \panel{page 46, Ada Lovelace and Grace Hopper}
  \panelslide{46}{5}{And many of those whose service jobs will be replaced belong to an already vulnerable group: minority women.}
\end{frame}
\begin{frame}{Feminism}
  \panel{page 46, Ada Lovelace and Grace Hopper}
  \panelslide{46}{6}{We need (different!) women to participate in the development and use of AI. Otherwise we'll just reproduce patriarchal structures.}
\end{frame}

% Reads AI through Weber's judge as automaton and Marx's worker as appendage
% of the machine, then names the ask. Credited on the page to Timo Daum.
\setpagethesis{The paradigm is not production and sale. It is collection and use.}
\begin{frame}{Digital Capitalism}
  \panel{page 47, Weber's judge and Marx's worker}
  \panelslide{47}{1}{Algorithms have always been an integral part of our modern capitalist societies.}
\end{frame}
\begin{frame}{Digital Capitalism}
  \panel{page 47, Weber's judge and Marx's worker}
  \panelslide{47}{2}{The idea of objectiveness, general abstract rules, quantification, de-personalization and the concept of equality are characteristics of modernity.}
\end{frame}
\begin{frame}{Digital Capitalism}
  \panel{page 47, Weber's judge and Marx's worker}
  \panelslide{47}{3}{Max Weber, e.g., pictures the judge as an automation who processes concrete cases (data) based on an algorithm, namely the law.}
\end{frame}
\begin{frame}{Digital Capitalism}
  \panel{page 47, Weber's judge and Marx's worker}
  \panelslide{47}{4}{Marx' workers are an appendage of the machine, alienated from the spiritual potencies of the capitalist working process.}
\end{frame}
\begin{frame}{Digital Capitalism}
  \panel{page 47, Weber's judge and Marx's worker}
  \panelslide{47}{5}{Digital capitalism is the ground on which AI thrives. The central paradigm is -- instead of production and sale -- the collection and use of data.}
\end{frame}
\begin{frame}{Digital Capitalism}
  \panel{page 47, Weber's judge and Marx's worker}
  \panelslide{47}{6}{Since today's AI is largely incomprehensible to us, we need to regulate and shape AI as a public service and snatch the omnipotence over AI from the few dominant, digital companies.}
\end{frame}

% After Virginia Dignum. Three pillars, one per definition panel, and panel 6
% refuses to let the acronym stand in for the work.
\setpagethesis{The acronym is the easy part. The governance it asks for does not exist.}
\begin{frame}{ART: Accountability, Responsibility, Transparency}
  \panel{page 48, the three equal pillars}
  \panelslide{48}{1}{Throughout the world, we share a growing awareness: We need to make AI safe, beneficial and fair for us.}
\end{frame}
\begin{frame}{ART: Accountability, Responsibility, Transparency}
  \panel{page 48, the three equal pillars}
  \panelslide{48}{2}{This requires the participation and commitment from many of us. It means training, regulation and awareness in 3 equal pillars:}
\end{frame}
\begin{frame}{ART: Accountability, Responsibility, Transparency}
  \panel{page 48, the three equal pillars}
  \panelslide{48}{3}{Accountability is the concept that AI should be held responsible for the results of its algorithms.}
\end{frame}
\begin{frame}{ART: Accountability, Responsibility, Transparency}
  \panel{page 48, the three equal pillars}
  \panelslide{48}{4}{This includes our responsibility to develop a framework for AI that incorporates our values.}
\end{frame}
\begin{frame}{ART: Accountability, Responsibility, Transparency}
  \panel{page 48, the three equal pillars}
  \panelslide{48}{5}{Transparency refers to the need to describe, inspect and reproduce the AI algorithms and results, and to manage the data used, in a fair way.}
\end{frame}
\begin{frame}{ART: Accountability, Responsibility, Transparency}
  \panel{page 48, the three equal pillars}
  \panelslide{48}{6}{But let's not fool ourselves. For this ART of AI, we need a new and more ambitious form of governance.}
\end{frame}

% Closes the Outlook on numeracy: none of the demands above can be argued by
% people whose brains signal pain at a number.
\setpagethesis{Who in this room went quiet when the first equation appeared?}
\begin{frame}{No Fear of Numbers}
  \panel{page 49, 1908 against MCMVIII}
  \panelslide{49}{1}{In order to understand how AI could improve our life, we propose a real revolution:}
\end{frame}
\begin{frame}{No Fear of Numbers}
  \panel{page 49, 1908 against MCMVIII}
  \panelslide{49}{2}{Healing the public shaming in our individual math-biographies that caused fear.}
\end{frame}
\begin{frame}{No Fear of Numbers}
  \panel{page 49, 1908 against MCMVIII}
  \panelslide{49}{3}{Many of our brains signal `pain' when we read numbers, think or talk about them.}
\end{frame}
\begin{frame}{No Fear of Numbers}
  \panel{page 49, 1908 against MCMVIII}
  \panelslide{49}{4}{Not only does this deprive life. Zero, e.g., is a fascinating tool allowing us to write `1908' instead of MCMVIII: 1,000+(1,000-100)+5+1+1+1.}
\end{frame}
\begin{frame}{No Fear of Numbers}
  \panel{page 49, 1908 against MCMVIII}
  \panelslide{49}{5}{Zero is also part of the binary code, fundamental to todays' computers and their language. Fundamental to AI, therefore.}
\end{frame}
\begin{frame}{No Fear of Numbers}
  \panel{page 49, 1908 against MCMVIII}
  \panelslide{49}{6}{With improved `math-esteem', more of us could and should discuss chances and risks of AI in our lives.}
\end{frame}
